\section{Pricing (Jake)}\label{sec:pricing}

Pricing is where real competition happens: the low-cost firm cuts its price, followers
match, and cuts stop at cost (Shaikh 2016, Table 7.3). Version 2 has no anchor and no war
trigger: a firm cuts only when it is cheaper than a rival at the lowest price, so price wars
come from cost differences, and technical change is what renews them. The group's
leader-follow rule, in which every firm cuts whenever capacity is not short, is the other
end of the price-war treatment; the two rules are the model's most sensitive switch, and
the sensitivity analysis decides whether one is fixed or both are kept as a treatment
(\code{docs/design/review-v2.md}).

\paragraph{Rule} (\code{PRICING\_RULE = "real\_competition"}). With
$R_{i,t} = \min_{j\neq i} p_{j,t-1}$ the lowest rival price,
$uc_i = \min_{v}(ucs_v + \delta\kappa_v)$ the unit cost of firm $i$'s best plant, $E_{t-1}$
last period's demand above normal capacity and $D_{t-1}$ last period's demand,
\[
p_{i,t} =
\begin{cases}
\min\{\, p_{j,t-1} : p_{j,t-1} > p_{i,t-1} \,\} & \text{if } E_{t-1} > 0
  \text{ and some rival priced higher}, \\
(1 + \eta\, E_{t-1}/D_{t-1})\, p_{i,t-1} & \text{if } E_{t-1} > 0
  \text{ and no rival priced higher}, \\
\min\{\, p_{i,t-1},\ (1-c)\,R_{i,t} \,\} & \text{if } u_{i,t-1} < 1 \text{ and }
  uc_i < \min\{(1-c)R_{i,t},\ \max_{j:\,p_{j,t-1}=R_{i,t}} uc_j\}, \\
\max\{\, \min\{p_{i,t-1}, R_{i,t}\},\ uc_i \,\} & \text{otherwise,}
\end{cases}
\]
with $c = 0.02$ and $\eta = 0.5$. In words: under excess demand, move up to the next higher
price, and the firm at the top raises it; with spare capacity and a cost edge over a rival
at the lowest price, undercut it; otherwise follow the lowest price, never below one's own
unit cost. The leader is whoever posts the lowest price: it is not assigned, and the
lowest-cost firm leads in 187 of 200 periods with the version 1 plants and accelerator
(seeds 1--10, $\rbar = 0.16$), in 58 under the version 2 defaults.

\paragraph{Why each detail.} $R$ excludes the firm itself, or a lone leader would undercut
itself every period. A cut never raises the firm's price. The edge is over \emph{a} rival at
$R$, so an identical twin at $R$ does not block the cut. The top-priced firm must be able to
raise, or prices could only fall. Excess demand is measured against normal capacity
(\code{EXCESS\_DEMAND = "normal\_capacity"}): against all output with overtime
(\code{"unserved"}), the 20\,\% overtime margin absorbs every shortage and prices end at
C's unit cost, 78.5. The floor is the firm's own unit cost; \code{PRICE\_FLOOR =
"operating\_cost"} uses the version 1 floor ($ucs$, Shaikh 2016, p.~316), with which
followers match the leader all the way down.

\paragraph{What moves the price in version 2.} Under the version 2 defaults the price does
not get stuck: retained earnings keep buying plants, idle scrapping keeps capacity near
demand, and each adopted design gives its owner a cost edge that starts a cut (a cut in 126
of 200 periods at seed 42). Shortages come from scrapping, not from under-investment. With
the version 1 plants and accelerator ($\rbar = 0.16$) the war of t = 0--19 drove prices to
78.5, a shortage at t = 20 lifted them to 89.9, and then nothing moved: capacity was rebuilt
to last sales and no firm could cut. The scratch test with profit-sensitive investment,
desired capacity $x_{t-1}(1+\beta(r_i-\rbar)/\rbar)$ with $\beta = 0.25$, under which the
price gravitates around 99 and the profit rate around 15.7\,\%
(\code{docs/design/choices.md}), remains the proposal for the accelerator arm of the
finance treatment.

\paragraph{Version 1} (\code{PRICING\_RULE = "anchor"}), kept for comparison. Firms move
towards an anchor, stay within a band of the leader, and cut collectively after a run of low
utilisation:
\begin{gather*}
p^N_i = uc_i + \rbar\,\kappa_i \quad\text{(capacity-weighted over the firm's plants)}, \qquad
p_{i,t} = \max\{\, ucs_i,\ \tilde p_{i,t} \,\}, \\
\tilde p_{i,t} =
\begin{cases}
(1-c)\, p_{i,t-1} & \text{if } u_{i,t-k} < u^{war},\ k = 1,\dots,T_{war}, \\
\min\{\, p_{i,t-1} + \lambda\,(p^{N}_i - p_{i,t-1}),\ (1+b)\, p^{lead}_{t-1} \,\} & \text{otherwise,}
\end{cases}
\end{gather*}
with $u^{war} = 0.95$, $T_{war} = 2$, $b = 0.025$, $\lambda = 0.25$.
\code{PRICING\_ANCHOR = "markup"} uses $(1+\mu)uc_i$; \code{"leader"} gives every firm the
target-return price of last period's leader. Over seeds 1--10 the final share of Shaikh
firms is 0.12 under each firm's own target return, 0.82 under the markup and 0.51 under the
leader's anchor: the anchor decides the winner.

\paragraph{Alternative: follow and undercut the leader} (\code{PRICING\_RULE =
"leader\_follow"}, proposed from the group's notebook). With $p^{lead}_{t-1}$ the lowest price
posted last period and a floor at a markup on average unit cost,
\begin{gather*}
\underline p_{i} = (1+m)\,\overline{uc}_i \quad\text{(capacity-weighted over the firm's plants)}, \\
p_{i,t} =
\begin{cases}
\min\{(1+z)\,p_{i,t-1},\ \bar m\,\underline p_{i}\} & \text{if } i \text{ is the only firm}, \\
(1+z)\,p^{lead}_{t-1} & \text{after excess demand}, \\
\max\{\underline p_{i},\ (1-c)\,p^{lead}_{t-1}\} & \text{otherwise},
\end{cases}
\end{gather*}
with $c = 0.02$, $m = 0.20$, $z = 0.03$, $\bar m = 3$ (\code{PRICE\_CUT}, \code{FLOOR\_MARKUP},
\code{LEADER\_RISE}, \code{MONOPOLY\_CAP}) and excess demand measured as in the default rule.
Every firm cuts whenever capacity is not short, so prices fall each period to the floors, and
the firm with the lowest floor sets the price. Unlike the notebook, $c$, $m$ and $z$ are the
same for every firm rather than drawn at random for each, so identical firms stay identical
and the label swap holds. Over seeds 1--10 under the version 2 defaults the leader is the
lowest-cost firm in 195 of 200 periods (58 under the default rule) and the final share of
Shaikh firms is 0.55 (0.81). Two points of the rule are open for the sensitivity analysis:
the floor, a 20\,\% markup on the firm's average unit cost under which no plant ever sells
below its unit cost and the Okishio hurdle never reaches zero (the version 2 default stops
cuts at the best plant's unit cost, version 1 at operating cost); and the shortage branch,
under which every firm posts $(1+z)\,p^{lead}_{t-1}$, so that a firm priced above that
lowers its price in a shortage.

\paragraph{Tests} (\code{tests/test\_pricing.py}): the rule's branches one by one (cut with an
edge and spare capacity; no cut without either, nor against an identical twin; a cut never
raises a price; the leader does not undercut itself; follow down to the floor and no further;
excess demand moves prices up); no price below the floor and no rise without excess demand
over 2\,000 random cases; identical firms with spare capacity never start a war; on the test
industry D2 cuts first, the others follow to their unit costs, and excess demand raises every
price; with the operating-cost floor the war ends on one price; the label swap holds under each
switch. Version 1 keeps its tests (1)--(4) and the leader-anchor tests. The leader-follow
rule: its branches one by one, prices falling to $1.2\,uc$ on the test industry, identical
firms posting identical prices, and the label swap.

\subsection*{Code}
\eqfile{anchor_price.py}
\eqfile{price_floor.py}
\eqfile{price_war.py}
\eqfile{real_competition_price.py}
\eqfile{leader_follow_price.py}
\eqfile{posted_price.py}
\schedulerpart{pricing}{\code{model/scheduler.py}: step 1}
